\begin{figure*}
\begin{minipage}[b]{0.30\textwidth}
{\small\begin{verbatim}
int
f(bool cond, int z) {
  int x, y;
  if (cond) {
    x = 3 * z;
    y = z;
  } else {
    x = 2 * z;
    y = 2 * z;
  }
  return x + y;
}
\end{verbatim}}
\centering (a)
\end{minipage}
%
\begin{minipage}[b]{0.40\textwidth}
{\small\begin{verbatim}
define i32 @f(i1 %0, i32 %1) {
  br i1 %0, label %3, label %5

label %3:
  %4 = mul nsw i32 %1, 3
  br label %8

label %5:
  %6 = shl nsw i32 %1, 1
  %7 = shl nsw i32 %1, 1
  br label %8

label %8:
  %.07 = phi i32 [ %4, %3 ], [ %6, %5 ]
  %.0 = phi i32 [ %1, %3 ], [ %7, %5 ]
  %9 = add nsw i32 %.07, %.0
  ret i32 %9
}
\end{verbatim}}
\centering (b)
\end{minipage}
%
\begin{minipage}[b]{0.33\textwidth}
{\small\begin{verbatim}
%0 = block 2
%1:i32 = var
%2:i32 = shlnsw %1, 1:i32
%3:i32 = phi %0, %1, %2
%4:i32 = mulnsw 3:i32, %1
%5:i32 = phi %0, %4, %2
%6:i32 = addnsw %3, %5
infer %6
\end{verbatim}}
$\Rightarrow$
{\small\begin{verbatim}
%7:i32 = shl %1, 2:i32
result %7
\end{verbatim}}
\centering (c)
\end{minipage}
\caption{A C++ function that can be optimized to ``\texttt{return z <<
    2}'' (a), its representation in LLVM~IR (b), and a corresponding
  \tool{} LHS and RHS (c).  The first line of the LHS defines a block
  value with two predecessors; each phi node attached to this block
  will have two regular value arguments.  Then, the first argument to
  each phi node, \texttt{\%0}, establishes that they are correlated:
  since they come from the same basic block, they must make the same
  choice. Without information about the correlation between the phi
  nodes, this optimization cannot be synthesized.}
\label{fig:phi}
\end{figure*}
